\section{Conclusions}\label{sec:conclusions}
\sectionstate{partial}{first prose draft 2026-11-18}
A forward model of the sampled strip signals, with the drift diffusion, the prompt footprint and the spreading along
the resistive strips as its only lateral terms and no fitted coupling between strips, reconstructs micro-TPC track
segments on a resistive-strip Micromegas without the bias of hit times. All of its constants are measured on the
chamber's own data, per chamber and per run condition. On five chambers on a cosmic bench it gives an angle
resolution of 0.9 to 1.1\textdegree{} on the two uniform chambers and 1.2 to 2.4\textdegree{} on the three with known
defects, where the hit line gives 2.0 to 2.6\textdegree{} even with its angle scale calibrated on the reference; at
a known angle of 25.4\textdegree{} in the SPS beam it reads $0.98\pm0.03$ of the stage tangent.

Fitting the same model to hit lists, one time and one amplitude per strip, removes the bias of the hit line and
halves its spread on the bench (1.2 to 1.3\textdegree{} against 1.0\textdegree{} for the waveforms), but leaves a
third of the beam planes more than 5\textdegree{} off: the model helps the hits and does not replace the samples.

The resolution is three times the limit set by transverse diffusion, 0.34\textdegree. The reference, the alignment
and the known event conditions do not account for the difference; it behaves as a fluctuation of the charge pattern
from event to event that is larger in the data than in our simulation, and is the main open point of the method.
